\begin{afterword}
\summaryhead

The post argues that the wonder of evolution is that it works at all. Evolution ends the regress of designers, who designed the designer, because it needs no cleverness: it is stupid and inefficient and works anyway.

The creationist picture of parts shaken in a box misdescribes evolution. Complex machinery evolves in small steps, or by putting old machinery to new uses, as gliding squirrels might one day become fliers. Parts that are useful alone can become mutually dependent, so that the finished machine breaks if any part is removed. The example is DNA and the ribosome, each useless without the other, which the post explains by RNA: it can both store information and do chemistry, though badly. Evolution does not claim to explain the first replicator, which was a crude accident.

Praising evolution beyond what science supports is inaccuracy, not loyalty to science. Evolution is not a designer for humans to imitate, and the post ends with Huxley's call to combat the cosmic process in ethics.

\responsehead

The post's main point is sound and worth making. Evolution is crude, and the reply to the watch-in-a-box objection is that complex parts are built in useful steps, not assembled at once. The post is also right that evolution does not claim to explain the origin of life. Its account of interlocking parts is the standard one, and an old one: Hermann Muller described such ``interlocking'' in 1918.

The worked example is stated more firmly than the evidence allows. The post says ``in this case we do know how it happened,'' and ends the same paragraph with ``Almost certainly.'' The hedge is the accurate one. RNA before DNA and protein has broad support, and the ribosome's RNA core is good evidence for it. But the paragraph gives only that first step, and how an RNA system became one based on DNA and proteins is still listed as an open problem. The next paragraph says, without a hedge, that RNA's ingredients ``would have arisen naturally'' on the early Earth. At the time this was an open problem: two years later the chemists who found a route to half of RNA's nucleotides wrote that it was ``far from obvious'' how they could form.

The warning against praising evolution too highly is aimed at an imagined speaker; the post quotes no one who makes the error. Its claim that ``Science has a very exact idea of the capabilities of evolution'' is supported the next day with formulas that are approximations for an idealized population, adequate for orders of magnitude.

The conclusion joins two claims. Evolution is not a guide to conduct, and Huxley's sentence, which is about ethics, supports that. Evolution is also a designer that humans should not imitate, and for this the post offers only the analogy of a modern bacterium imitating the first replicator. By then engineers were imitating evolution on computers with some success: evolved antennas flew on a NASA mission in 2006. The analogy fits evolution as it runs in nature, not every imitation of it.

In short: a sound point about evolution's crudeness, with a worked example stated more firmly than its own hedge, and a verdict against imitating evolution as a designer that is broader than the evidence.
\end{afterword}
